% \jh{we should have a figure for the system}
% \jh{this section should be expanded further}
% \jh{In the intro part we listed several principles of Agentboard. Intro is only a summary, typically we should recap and re-emphasize the principles and draw comparisons with previous works again, and cite Table 1 again (Table 1 is an important table, but we only briefly mention it in the intro across the paper). However, it seems we never talk about these principles and comparisons after intro, which is important. When we talk about the comparison with Agentbench, the observational law paper can be cited. }

\BenchName{} is a unified, open-source benchmark for evaluating LLM agents that adheres to five key principles: \textit{task diversity, multi-round interaction, partially-observable environments, fine-grained metrics, and analytical evaluation}, as shown in Table~\ref{table: dataset_comparison}. Our commitment to these principles manifests in three key areas:
\begin{itemize}[before=\vspace{-1mm}, leftmargin=10pt]
    \item \textit{Task Diversity and Uniformity}, 
    where we carefully curate nine diverse environments across four scenarios, ensuring they require multi-round interactions, are fully text-based and primarily partially-observable. This contrasts with many existing LLM agent benchmarks, which are often solvable in a single round, derived from fully-observable tasks like MMLU, or focus on a specific type of task, as demonstrated in Table~\ref{table: dataset_comparison}. 
    % These designs provided by our benchmark enable a comprehensive evaluation of agent capabilities, as shown in analysis by
    Compared to AgentBench~\citep{liu2023agentbench}, in particular, our choice of tasks makes \BenchName{} planning-heavier, where the results on \BenchName{} well-correlates with the scores on traditional reasoning and coding benchmarks, while the AgentBench scores strongly correlated with scores on the knowledge test MMLU, as illustrated in
    \citet{ruan2024observational}. 
    % demonstrated that results on \BenchName{} correlates well with standard reasoning and coding benchmarks assessing reasoning and coding—two cornerstones of agent abilities. 
    % This contrasts with benchmarks like AgentBench that tend to focus predominantly on factual knowledge, underscoring \BenchName{}'s broader applicability in assessing general LLM agent capabilities. 
    We elaborate on the choice of tasks and their adaptation in \S\ref{section: task decomposition}.
    
    \item \textit{Fine-grained Progress Rate}, where \BenchName{} is the first to propose a fine-grained progress rate metric tracking the intermediate progress of different agents. This metric distinguishes our benchmark in tracking minimal improvement in LLM agent performances. Such a capability is crucial in current endeavors to develop stronger open-weight LLMs, providing detailed insights that are essential for incremental advancements in agent capabilities. We provide detailed introduction for this metric and its annotation process in \S\ref{sec:fine-grained progress rate}.
    
    \item \textit{Comprehensive Analysis}, where \BenchName{} is the first LLM Agent benchmark to expand metrics beyond mere success rate and scores to include detailed analyses.
    As illustrated in Figure~\ref{fig:overall illustration}, such a comprehensive evaluation includes (1) fine-grained progress rates tracking different agents, (2) grounding accuracy, (3) performance breakdown for hard and easy examples, (4) long-range interactions, (5) analyses of performance across various sub-skills, and (6) trajectory with friendly visualization. 
    We elaborate these analyses in our experiments at \S \ref{sec:exp}.
    Additionally, \BenchName{} provides an web interface\footnote{\label{wandb link}An example of the panel is public \href{https://wandb.ai/agentboard/llm-agent-eval-gpt-4-all}{here}.} through Wandb dashboard that offers interactive visualizations of these analyses during evaluation. We perform a case study on the panel in \S \ref{sec: wandb panel}.
\end{itemize}


% \BenchName{} is designed around the core principles of uniformity and user-friendliness. Beyond a text-based agent evaluation benchmark, we develop \BenchName{} aiming for a readily accessible, open-source evaluation framework that facilitates diverse analysis and easy customization for different models, agents, and environments, all within a unified format. 
% This commitment to uniformity manifests in three key areas:
% \begin{itemize}[before=\vspace{-3mm}]
%     \item \textit{Interface Uniformity}, where we present a consistent interface implementation across various environments, model deployments, and prompt types; 
%     \item \textit{Observation and Action Space Uniformity}, ensuring datasets are uniformly constructed for text-based interactions with unified metrics;
%     \item \textit{Analysis Uniformity}, which expands metrics beyond mere success rate and scores to include detailed analyses.
%     As illustrated in Figure~\ref{fig:overall illustration}, such a comprehensive evaluation includes assessment on (1) fine-grained progress rates tracking different agents, (2) grounding accuracy, (3) performance breakdown for hard and easy examples, (4) long-range interactions, (5) analyses of performance across various sub-skills, and (6) trajectory with friendly visualization. 
%     We elaborate these analyses in our experiments at \S \ref{sec:exp}.
%     Additionally, \BenchName{} provides an web interface through Wandb dashboard that offers interactive visualizations of these analyses during evaluation.
%     % \footnote{\label{wandb link}An example of the panel is public \href{https://wandb.ai/agentboard/llm-agent-eval-gpt-4-all}{here}.} 
%     We perform a case study on the panel in \S \ref{sec: wandb panel}.
%     % Figure \ref{fig:overall illustration} visually represents the extensive range of analysis incorporated within our evaluation system.
% \end{itemize} 



\subsection{A Unified Multi-Round Reflex Agent}
\label{section: agent structure}

\paragraph{Preliminaries:}
An LLM agent receives textual world descriptions, chooses a text action, and gets feedback detailing state changes and any action errors. Interaction with these environments can be modeled as a special case of Partially Observable Markov Decision Processes (POMDPs) defined by tuple $\langle g, \mathcal{S}, \mathcal{A}, \mathcal{O}, \mathcal{T}\rangle$, with goal $g$, state space $\mathcal{S}$, valid actions space $\mathcal{A}$, observation space (including environment feedback) $\mathcal{O}$, transition function $\mathcal{T}: \mathcal{S} \times \mathcal{A} \rightarrow \mathcal{S}$.
An agent with policy $\pi$ makes prediction at time step $t$ based on goal $g$ and memory $m_t = \{o_j, a_j, o_{j+1}, a_{j+1}, \dots o_t\}, 0 \leq j < t$, which is a sequence of actions and observations. This trajectory of the agent $\tau=[s_0,a_0,s_1,a_1,\dots s_t]$ is formulated by policy and environmental state transitions, such as
\small
\begin{equation}
    p_{\pi}(\tau) = p(s_0)\prod_{t=0}^T\pi(a_t|g, s_t, m_t)\mathcal{T}(s_{t+1}|s_t, a_t)
\end{equation}

% In \BenchName{}, environments are deterministic, with the agent's trajectory driven only by the language model's policy. 

% The agent's structure is detailed in \S \ref{section: agent structure}, and interactive environment designs in Section \ref{section: task decomposition}.

\begin{figure*}[!t]
    \centering    
    \includegraphics[width=\textwidth]{figures/agent_prompt.pdf}
    \label{dataset}
    \vspace{-3mm}
    \caption{(Left) A structural overview of the reflex agent, which iteratively interacts with the environment and makes next step predictions based on the goal and history. (Right) An example of a prompt for our reflex agent.}
    \label{fig: agent structure}
\end{figure*}

% \subsection{A Unified Reflex Agent\label{section: agent structure}}


\paragraph{The Unified Framework:} 
\BenchName{} unifies all tasks around a general framework where the agent receives observations $o_t$ and performs actions $a_t$, causing deterministic state transitions $\mathcal{T}: (s_{t}, a_t) \rightarrow s_{t+1}$ based on real-world dynamics. A feedback function $f$ is also defined in the environment to derive feedback from each interaction round $o_t = f(s_t, a_t)$. This feedback includes: (1) list all valid actions when the agent uses help actions such as \textit{check valid actions}; (2) execute valid action $a_t$ and return a description of the changed stete $s_{t+1}$; (3) issue errors when the agent performs an action outside of the action space. 


% \BenchName{} unifies all tasks around a general framework where the agent iteratively takes observation/feedback $o_t$ from the environment and outputs actions $a_t$. The action then acts upon the environment and transitions into a new state $s_{t+1}$. The transition function $\mathcal{T}: (s_{t}, a_t) \rightarrow s_{t+1} $ of all task environments are deterministic and defined in the environments according to real-world rules. 
% A feedback function $f$ is also defined in the environment to derive feedback from each interaction $o_t = f(s_t, a_t)$. This feedback handles three kinds of scenarios: (1) When the agent uses help actions such as \textit{check valid actions}, the environment returns all valid actions; (2) When the agent generates executable function, the environment returns a description of the changed state $s_{t+1}$ by executing the action; (3) When the agent performs an action not in action space $\mathcal{A}$, the environment returns an error message to help with debugging and troubleshooting.

% We integrate various task-solving processes into a unified reflex agent framework that makes decisions based on past experiences, 
We aim to use a simplistic agent framework to showcase LLM basic agentic abilities. As shown in Figure \ref{fig: agent structure}, our agent makes decisions based on its memory of past perceptions,
similar to how humans learn from experience and adapt. The implementation of the reflex agent assessed in this paper adopts an act-only prompting strategy in line with recent studies \citep{liu2023bolaa, zhou2023webarena, xu2023tool}, detailed in the right part of Figure \ref{fig: agent structure}, while other prompting strategies can be easily incorporated into our open-source framework. 
% A task is introduced with an instruction and an example, followed by a clear goal and an interaction memory formatted with \textit{``Observation: ''} and \textit{``Action: ''} prompts to guide the LLM. 
% The prompt template and content are kept consistent across various LLMs, with occasional minor adjustments to fit specific model requirements. 
Also, LLM agents tend to struggle with limited context lengths in long interactions, failing to retain full history. Following the \textit{``sliding window''} method from LangChain~\citep{Chase_LangChain_2022}, we focus on recent, more impactful interactions \citep{puterman1990markov} within context constraints. This differs from previous practices that stop the agent when context limits are surpassed \citep{liu2023agentbench, wang2023mint}, allowing for extended, intricate interactions in our approach. We provide ablation results such as ReAct prompting~\citep{yao2022react} and other long-context processing techniques to justify our framework in Appendix \ref{sec: ablation study}.
% Also, LLM agents often struggle with limited context lengths, unable to capture the full history in long interactions. Following the \textit{``sliding window''} method from LangChain~\citep{Chase_LangChain_2022}, we prioritize recent interactions, deemed more critical for future actions~\citep{puterman1990markov}, within the context limit. This contrasts with earlier methods that halt the agent when exceeding context thresholds~\citep{liu2023agentbench, wang2023mint}, whereas our approach facilitates longer, more complex interactions.

% In our design, we unify problem-solving of various tasks under a unified reflex agent framework. Currently, we adopt our agent structure as the classic goal-based reflex agent in \citet{russell2005ai}. As shown in Figure \ref{fig: agent structure}, the agent selects actions on the basis of the perception history, i.e., memory. For example, if the agent wants to pick up a thermometer from a cupboard, and the environment informs it that this action is not executable, as there is no thermometer in the cupboard. Therefore, the agent adapts to the feedback and look for the thermometer elsewhere, such as, the agent establish a connection between the last feedback from the environment to the next action. Humans also have the ability to make such connections, enabling us to learn from trial and errors, and adapt to our surroundings \citep{piaget2005psychology}. We also expect LLM agents to have such ability. 

% In this paper, we implement the reflex agent based on an act-only prompting approach following~\citep{liu2023bolaa, zhou2023webarena, xu2023tool}, while other prompting approaches such as ReAct~\citep{yao2022react} can be applied as well. As shown in the right part of Figure~\ref{fig: agent structure}, first, an overall instruction of the task is given, followed by a single in-context learning example. Then, the goal is clearly stated such that the model would perform actions towards finishing this problem. The interaction memory is provided in the format of stating \textit{``Observation: ''} and \textit{``Action: ''} so that the model could understand the impact of the action on the environment, as well as how the world evolves. 
% % These items are combined together as a prompt to the LLM. 
% The uniformity of the prompt template and contents is maintained across diverse LLMs, though minor modifications to the prompt format are occasionally necessary to accommodate the requirements of different models. 

% A common challenge for LLM agents is that they have a limited context length that often fails to record all interactive history across long-range interactions. We follow  the ``sliding window'' approach in LangChain~\citep{Chase_LangChain_2022}, maintaining only the most recent interactions within the maximum context length. This approach is based on the assumption that recent interaction records, which describe the current state of the agent, are more important than earlier interactions for future actions~\citep{puterman1990markov}. Unlike previous approaches such as those described in the works of ~\citet{liu2023agentbench, wang2023mint} which stop the agent when the context length exceeds a certain threshold, our approach allows for benchmarking longer interactions and enables the agents to handle more complex tasks.



% \subsection{Multi-turn Interaction Framework}
\subsection{Fine-grained Progress Rate \label{sec:fine-grained progress rate}}

Recent studies highlight the predominant use of success rate as the main metric for agent evaluation, which fails to capture the nuances of partial task completion by language model agents~\citep{liu2023agentbench, li2023api}. This approach does not differentiate between near-complete tasks and minimal task execution, treating both as equivalent failures. Alternative metrics like reward scores are available but lack standardization~\citep{chevalier2018babyai, wang2022scienceworld}. To mitigate this issue, we introduce a \emph{progress rate} metric to accurately reflect LM agents' goal attainment at various stages.

% \renewcommand{\arraystretch}{1.2} % 1.2倍行高
\begin{table*}[t!]
\centering
\scriptsize
\caption{Examples of goals for the 4 task categories in \BenchName{}, along with a sampled step of the trajectory and progress rate.
The trajectory is generated by GPT-4.
Some lengthy observations are omitted with ``$\ldots$'' for brevity. The task name in the table uses an abbreviation, the full name can be found in \S \ref{section: task decomposition}.
}
\vspace{-2mm}
\resizebox{\textwidth}{!}{
\begin{tabular}{p{0.3cm}|p{12.50cm}}
\toprule
\multicolumn{1}{c|}{\textbf{Task}}  & \multicolumn{1}{c}{\textbf{Goal \& Trajectory}} \\
\midrule
\midrule
\multirow{2}[5]{*}{\textbf{ALF}} 
&
\textbf{Goal: put a clean egg in microwave.}
% id: 23
\\
&
\stepbox{Step 02}
\hspace{0.18cm}\tcv{Action:} open fridge 1
\ \tco{Observation:} You open the fridge 1. The fridge 1 is open. In it, you see a apple 2, a egg 1, a lettuce 1, a pan 2, a plate 1, and a tomato 1.
\hfill \stepbox{Progress Rate: 0.00 $\rightarrow$ 0.25}
% 218, 238, 253
% 183, 220, 252  0.71, 0.86, 0.99
% \\
% \midrule
% \midrule
% \multirow{2}[5]{*}{\textbf{SW}} 
% &
% \textbf{Goal: Your task is to boil apple juice until it reaches a gaseous state.}
% % id: 36
% \\
% &
% \stepbox{Step 25}
% \hspace{0.05cm}\tcv{Action:} use thermometer in inventory on apple juice in metal pot
% \ \tco{Observation:} the thermometer measures a temperature of 23 degrees celsius
% \hfill \stepbox{Progress Rate: 0.75 $\rightarrow$ 0.88}
% \\
% \midrule
% \midrule
% \multirow{2}[5]{*}{\textbf{BA}} 
% &
% \textbf{Goal: pick up the box}
% % id: 59
% \\
% &
% \stepbox{Step 07}
% \hspace{0.05cm}\tcv{Action:} toggle and go through green closed door 1
% \quad\tco{Observation:}  In front of you in this room, you can see several objects: There is a red box 1 $\ldots$ 
% \hfill \stepbox{Progress Rate: 0.33 $\rightarrow$ 0.66}
\\
\midrule
\midrule
\multirow{2}[5]{*}{\textbf{JC}} 
&
\textbf{Goal: Get out of the house. Then escape the city without getting caught via driving.}
% id: 0
\\
&
\stepbox{Step 29}
\hspace{0.18cm}\tcv{Action:} take a shower
\ \tco{Observation:} You step into the shower, turn on the water, and within a few moments you feel like a new man. But no time to dawdle - you hop back out again and dry off in record time $\ldots$ \hfill \stepbox{Progress Rate: 0.43 $\rightarrow$ 0.57}
\\
\midrule
\midrule
% \multirow{2}[5]{*}{\textbf{PL}}
% &
% \textbf{Goal: The goal is to satisfy the following conditions: Wheel r1 is inflated., r1 is on the-hub1., w1 is in boot.}
% % id: 50
% \\
% &
% \stepbox{Step 08}
% \hspace{0.18cm}\tcv{Action:} Put-away w1 in boot.
% \ \tco{Observation:} W1 is in boot.
% \hfill \stepbox{Progress Rate: 0.00 $\rightarrow$ 0.33}
% \\
% \midrule
% \midrule
% \multirow{2}[5]{*}{\textbf{WS}} 
% &
% \textbf{Goal: i want a noise cancelling cosycost usb microphone, and price lower than 70.00 dollars}
% % id: 1
% \\
% &
% \stepbox{Step 02}
% \hspace{0.18cm}\tcv{Action:} search[noise cancelling cosycost usb microphone]
% \ \tco{Observation:} WEB PAGE: \{ [Back to Search] Page 1 (Total results: 15) [Next $>$] $\ldots$
% \hfill \stepbox{Progress Rate: 0.00 $\rightarrow$ 0.33}
% \\
% \midrule
% \midrule
\multirow{2}[5]{*}{\textbf{WA}} 
&
\textbf{Goal: Display the list of issues in the kkroening/ffmpeg-python repository that have labels related to questions}
% id: 1
\\
&
\stepbox{Step 05}
\hspace{0.18cm}\tcv{Action:} click [5398]
\ \tco{Observation:} Tab 0 (current): Issues · Karl Kroening / ffmpeg-python · GitLab [6573] RootWebArea `Issues · Karl Kroening / ffmpeg-python · GitLab' focused: True [6620] link $\ldots$
\hfill \stepbox{Progress Rate: 0.25 $\rightarrow$ 0.50}
% \\
% \midrule
% \midrule
% \multirow{2}[5]{*}{\textbf{TQ}} 
% &
% \textbf{Goal: Avatar versus Forrest Gump, which has a higher rating?}
% % id: 22
% \\
% &
% \stepbox{Step 04}
% \hspace{0.18cm}\tcv{Action:} get\_movie\_details with Action Input: \{``movie\_id": ``13"\}
% \ \tco{Observation:} \{`title': `Forrest Gump', `vote\_average': 8.476, `budget': 55000000, $\ldots$ 
% \hfill \stepbox{Progress Rate: 0.60 $\rightarrow$ 0.80}
\\
\midrule
\midrule
\multirow{2}[7]{*}{\textbf{TO}} 
&
\textbf{Goal: In ``Sheet17'', calculate and complete the ``Profit'' of the products in the table based on the sales information of the products. And then, sort the table in descending order by ``Profit''.}
% id: 36
\\
&
\stepbox{Step 07}
\hspace{0.1cm}\tcv{Action:} update\_cell\_by\_formula with Action Input: \{``operator'': ``PRODUCT'',  `` start\_position'': ``C8'', ``end\_position'': ``D8'', ``result\_position'': ``E8''\}
\ \tco{Observation:}
[[`Product', `Category',
%`Price',
%`Sold out', 
% `Profit'], [`banana', `Fruits', 
$\ldots$ 
\hfill \stepbox{Progress Rate: 0.47 $\rightarrow$ 0.49}
\\
\bottomrule
\end{tabular}}

\label{tab:goal_trajectory}
% \vspace{-5mm}
\end{table*}

% Current benchmarks for evaluating agents, as outlined in recent studies~\citep{liu2023agentbench, li2023api}, predominantly focus on the final success rate as the primary performance metric. Nevertheless, the complexity inherent in tasks designed for language model agents often results in these agents achieving only partial fulfillment of the intended goals, which is not accurately captured by solely considering the success rate,
% The difference between an agent that almost completes a task and one that only manages to perform a few steps is significant, yet success rates treat these outcomes as equivalent failures. 
% Several environments offer fine-grained metrics, including reward scores~\citep{chevalier2018babyai, wang2022scienceworld, hausknecht2020interactive}. However, these reward scores lack unified definition and vary in scales, making it difficult to assess overall task progress across multiple environments. To address these issues, we propose a new metric called \emph{progress rate} to checkpoint the current progress of LLM agents in accomplishing goals.

In each round of interaction, a progress rate, denoted as $r_t$, is assigned to evaluate the agent's advancement towards the goal state $g$. As the agent moves through the states $\mathbf{s_t}=[s_0, \dots, s_t]$, we assess its progress using a matching score $f(\cdot,g)\rightarrow [0,1]$ that quantifies the similarity between the current state and the goal state. The initial value of $r_t$ is set to 0, indicating no progress. The progress rate $r_t$ reflects the highest matching score achieved, reaching 1 when the task is completed. The progress rate is formulated as below: 
\begin{equation}
r_t = \left\{
\begin{aligned}
    &r_t^{\text{match}} = \max_{i, 0\leq i \leq t} f(s_i, g), \quad \, \, \text{if } f(\cdot, g)  \text{ is continuous} \\
    &r_{t}^{\text{subgoal}} = \max_{i, 0\leq i \leq t}\left(\frac{1}{K}\sum_{k=1}^Kf(s_i,g_k)\right), \quad \text{otherwise}
\end{aligned}
\right. 
\label{eq: progress rate}
\end{equation}
The function $f(\cdot, g)$ measures state similarity in tasks, such as comparing table states in manipulation activities. It works well for tasks with direct state comparisons but is less effective for tasks with ambiguous intermediate states, where progress is hard to measure. We mitigate this by introducing a discrete matching score to assess how closely intermediate states align with defined subgoals.
% The function $f(\cdot, g)$ measures task-specific state similarity; for instance, in table manipulation tasks it compares the current and final tables. This function is straightforward for direct manipulation tasks with comparable states. However, in our benchmark, most tasks, like ``clean an egg and put it in microwave", involve intermediate states that are not clearly indicative of progress toward the final goal. Assessing specific stages such as "open the fridge" is difficult. To tackle this, we employ human annotations alongside a subgoal-based distance metric. Goals are described textually, making it challenging to gauge exact advancement at a given time 
% $t$. To address this complexity, we introduce a discrete matching score that assesses how closely intermediate states align with defined subgoals.
We begin by decomposing the overall goal $g$ into a sequence of subgoals $\mathbf{g} = [g_1, \dots, g_K]$, with each subgoal leading into the next. The authors manually label each subgoal, which is then checked and adjusted through a rigorous process described in \S\ref{sec: annotation verification}.
Notably, we manually edit the problems for a simpler setup where each final goal aligns with a unique subgoal sequence, and this affects only 5\% of the original problems (a detailed descriptions for our adaptations are in Appendix \ref{sec: explanation for adaptation}). Note that while we maintain a unique subgoal sequence,  this allows for a diverse set of trajectories, e.g. taking detours when accomplishing the task.
% \footnote{Problems with more than one ground truth path are removed or rewritten.}
As an example, for task ``clean an egg and put it in microwave", the necessary subgoals would be ``open the fridge" $\rightarrow$ ``taking an egg from the fridge" $\rightarrow$ ``clean the egg with sinkbasin" $\rightarrow$ ``put the egg in the microwave". Each subgoal $g_i$ is associated with a labeled state that indicates its completion. 
To evaluate the match between an agent state and a subgoal, we employ a regular-expression-based matching function denoted as $f(\cdot, g_i)\rightarrow \{0,1\}$ and the progress rate as $r_t^{\text{subgoal}}$ in Equation \ref{eq: progress rate}.

We employ progress rate along with the commonly used success rate metric, which computes the proportion of tasks completed within $T$ interactions. 
% Success is achieved if the final progress rate $r_T$ equals 1, indicating that all subgoals ($r_t^{\text{subgoal}}$) or the goal state ($r_t^{\text{match}}$) have been reached.


% where the agent must finish $k$ consecutive subgoals to be awared a progress rate of $k/K$. Both $r_t^{\text{match}}$ and $r_{t}^{\text{subgoal}}$ record the maximum progress reached till step $t$, such as if the agent make mistakes and deviate further from the current goal, it would still be recognized for its previous progress. 
% Progress rate quantifies an LLM agent's efficiency in decomposing and executing subgoals. Discrete matching score, detailed in Section \ref{section: task decomposition}, serves as the progress rate for certain tasks, with human annotators confirming the subgoal sequence. The implementation of progress rate as a fine-grained metric could effectively facilitate the monitoring of preliminary progress in LLM agents, which is further discussed in \S \ref{section: fine-grained reward}.

% Progress rate measures the ability of the LLM agent to break down complex goals and perform the subgoals in an organized manner. 
% For tasks using the discrete matching score as progress rate, the sequence of subgoals is labeled and verified by human annotators, as detailed in \S \ref{section: task decomposition}. The implementation of progress rate as a fine-grained metric could effectively facilitate the monitoring of preliminary progress in LLM agents, which is further discussed in \S \ref{section: fine-grained reward}.

% \subsection{Fine-grained Progress Rate}
% \chang{I'm keeping this paragraph for comparison with the revised version. I believe we are causing confusion with matching score as a different way of calculating progress rate for several tasks. We can view the matching score as a continuous matching function between goal state and intermediate state, while our subgoal based score as a discrete matching metrics. I wrote a revised version of progress rate, that unifies our two ways of calculation.}

% Existing agent benchmarks ~\citep{liu2023agentbench, li2023api} primarily rely on the final success rate as a metric. However, due to the inherent difficulty of agent tasks for language models, agents typically accomplish only a subset of subgoals within these tasks. Consequently, using success rate alone as a metric fails to adequately measure the progress achieved by the agent. To address this, we propose a new metric called progress rate for LLM agents. This metric quantifies the percentage of manually annotated subgoals that the agent has successfully completed.

% Progress rate, denoted as $r_t$, is allocated for each round of interaction. Given that the goal is given in textual form, it is non-trivial to ascertain the precise phase of completion for the agent at a specific timestep $t$. \jh{I think we have discussed this, and when I was discussing with Shuyan, she also raised this point: there are actually different trajectories with different subgoals to reach the final goal -- at least this is true in realistic scenarios. This means theoretically the goal g may not be factorized as a sequence of subgoals (but multiple sequences)? Then how we annotate subgoals and how is it possible? Did we simplify the setting by testing on benchmarks with deterministic subgoals? Please discuss and clarify it here.}We first decompose the goal $g$ into a sequence of subgoals $G = [g_0, g_1, \dots, g_K]$, each preceding another, such that the task is only finished after step-by-step completion of all $K$ subgoals. For instance, if you want to clean yourself, you need to first go to the bathroom, take off your clothes, and then take the shower, \textit{going to the bathroom}, \textit{taking off clothes} and \textit{taking the shower} are necessary subgoals. For each subgoal $g_i$, we identify the corresponding state $s_{\alpha_i}$ that demonstrates that the subgoal is completed. The progress rate is defined as: 
% \begin{equation}  
% r_{t} = \argmax_{k, 0 \leq k \leq K}\prod_{i=0}^k \sigma[\sum_{j=0}^t\delta(s_{\alpha_i}, s_j)]/K, 
% \end{equation}  
% where $\sigma(\cdot)$ is 1 when the input is over 1, 0 otherwise, and $\delta(s_1, s_2)$ is an indicator function determining whether the two states are the same. The progress rate calculates the percentage of subgoals finished until this step, and the subgoals are required to finish in a sequential manner, such that $g_{i+1}$ is not allowed to finish before $g_{i}$. This score measures the ability of the LLM agent to break down complex goals and perform the subgoals in an organized manner. The deconstruction of objectives is labeled and verified by human annotators, and the computation of progress rates is incorporated within the environmental framework.
% The implementation of progress rate as a fine-grained metric could effectively facilitate the monitoring of preliminary progress in LLM agents. This is further discussed in Section \ref{section: fine-grained reward}.

% % The structure of the agent is then introduced in Section \ref{section: agent structure} and the task-specific designs of interactive environments are introduced in Section \ref{section: task decomposition}.



\section{\BenchName{} -- Task Composition\label{section: task decomposition}}
% composed of four types of scenarios, embodied game focus on navigation and searching, game focus on planning and strategies, web focus on understanding web and tool focus on grounding. 

% %\renewcommand{\baselinestretch}{1.0}
\begin{table*}[!t]
\tiny
\caption{Statistics of 9 environments in \BenchName. ``subgoal'' and ``match'' means 2 different implementations of progress rate $r^{\text{subgoal}}$ and $r^{\text{match}}$ respectively.  $^\dagger$Note that Sheet Environments in Tool-Operation are evaluated with $r^{\text{match}}$, while other sub-tasks are evaluated with $r^{\text{subgoal}}$. $^\ddag$For tasks without subgoal label, we state the average number of constraints to satisfy in the goal state, which is essentially the complexity of the problems. For context length, we report the number of tokens generated with Llama2 tokenizer. $^\dagger$ We divide problems into hard/easy based on the number of subgoals -- problems with a larger number of subgoals than cutoff are viewed as hard. }
\vspace{-3mm}
\centering
\resizebox{\textwidth}{!}{\begin{tabular}{@{}lccccccccc@{}}
\toprule
\linespread{0.2}

  \multirow{2}{*}{} & \multicolumn{3}{|c|}{\textbf{Embodied AI}} & \multicolumn{2}{c|}{\textbf{Game}} & \multicolumn{2}{c|}{\textbf{Web}} & \multicolumn{2}{c}{\textbf{Tool}}  \\

   & \multicolumn{1}{|c}{\textbf{ALF}} & \multicolumn{1}{c}{\textbf{SW}} & \multicolumn{1}{c|}{\textbf{BA}} & \multicolumn{1}{c}{\textbf{JC}} & \multicolumn{1}{c|}{\textbf{PL}} & \multicolumn{1}{c}{\textbf{WS}} & \multicolumn{1}{c|}{\textbf{WA}} & \multicolumn{1}{c}{\textbf{TQ}} & \multicolumn{1}{c}{\textbf{TO}}  \\  \midrule
\# Environment       & 134 & 90 & 112 &  20  & 60 & 251 & 245 & 60 & 40   \\
\# Turns          &  6 & 15 & 10 & 20 & 20   & 3 & 25 & 5  & 6   \\
Action Space  & 13 & 21 & 8 & 150 & 8   & 2 & 12 & 15 & 16  \\


\# Avg. Subgoals$^\ddag$     & 3 & 5 & 4 & 6 & 6 & 4   & 6 & 5 &  5  \\
Hard/Easy Cutoff$^\dagger$ & 3 & 3 & 3 & 4 & 6 & 1 & 4 &  4 & 4                \\
Context Length & 900 & 2800 & 1800 & 1500  & 2700   & 1200 & 15000 & 2100  &  4300  \\ 

Progress Rate             &    subgoal &  subgoal & subgoal  &    subgoal &  match & match & match & subgoal & subgoal/match $^\dagger$     \\

Success Rate &  \multicolumn{9}{c}{(Progress Rate == 1)} \\
 \bottomrule

\end{tabular}}
\
\label{tab: Statistics of 10 environments}
\end{table*}
% \jh{we should briefly talk about the principles of selecting these tasks here, aligning with the principles such as multi-round and partially-observable, and compare with previous works like agentbench and mint again}
\BenchName{} features four task scenarios: embodied, game, web, and tool. These tasks are selected for their diversity and relevance to everyday activities, offering broader scenario coverage than tool-using benchmarks like MINT~\citep{wang2023mint} and ToolEval~\citep{qin2023toolllm}. We specifically select tasks that require multi-round interactions in partially-observable environments, creating a realistic and challenging setting for agents. Examples of goals and trajectories are displayed in Table \ref{tab:goal_trajectory}, and task statistics are summarized in Table \ref{tab: Statistics of 10 environments}. Further details on the environments and annotation are provided in Appendix~\ref{appendix:detail environment} and~\ref{sec: Annotation of Subgoals}.

% \chang{!! write about four things 1. why is this environment important and what agentic abilities it evaluates. 2. how do we repurpose our environment (e.g. change to text, change to multiround) 3. What are the goals and action like: (see agentbench) 4. Metrics}

% \chang{Pls remove repetitive information about our fine-grained reward.}

\subsection{Environments}
%Assessing agents as action-performing robots is key to measure fundamental skills like space navigation, object search, and grounding. We use three varied tasks.




\textbf{\textit{Embodied} - AlfWorld (ALF)~\citep{shridhar2020alfworld}} are
household tasks that require agents to explore surroundings and perform commonsense tasks 
% Within \BenchName, we evaluate a model's ability to perform tasks in physical household settings, 
like ``put two soapbars in garbagecan". 
% AlfWorld is categorized into six types, comprising a total of 134 environments. 
This task uses subgoal-based progress rates (Appendix~\ref{sec: Annotation of alfworld}) and the original success rate of the environment as metrics. 
% Specifically, for each environment, we labeled N-1 observations that must be met to achieve the given goal, in the form of regular expressions, as N-1 subgoals. 
% For example, for the task: ``put a pencil on the desk'', one necessary observation we annotate is ``You pick up the pencil $\backslash$d +" (regular expression). It will match observations such as ``You pick up the pencil 1''. 
% If one observation matches the regular expression, the corresponding subgoal is then considered as achieved. 
% When the goal is completed, the environment returns a task success flag. 
% The final success flag combined with the N-1 annotated subgoals constitutes the set of N subgoals.

\textbf{\textit{Embodied} - ScienceWorld (SW)~\citep{wang2022scienceworld}}  is a challenging interactive text environment testing scientific commonsense, e.g.``measure the melting point of the orange juice".  The current subgoals provided by SW do not accurately reflect a language model's performance due to their sparsity and uneven weighting, as further explained in Appendix~\ref{sec: Re-annotation of Scienceworld}. To rectify this, we re-annotate the subgoals for calculating progress rate $r_t^{\text{subgoal}}$. We also re-annotate instructions on tool usage and rooms to explore to ensure the uniqueness of subgoal sequence for task completion.

%\paragraph{ScienceWorld (SW)~\citep{wang2022scienceworld}} 
%is a complex interactive text environment that poses a significant challenge to agents' scientific commonsense. %This environment requires agents to complete tasks such as ``measure the melting point of the orange juice".  
%While ScienceWorld has already provided subgoals, these subgoals pose challenges in reflecting the performance of a LLM due to the subgoal sparsity and inappropriate weighting of different subgoals as further explained in Appendix~\ref{sec: Re-annotation of Scienceworld}.
%navigate through 8 distinct functional rooms (e.g., workshop, kitchen) and utilize the tools to  
% In the original labels, subgoals are divided into ``sequential subgoals'' and ``unordered and optional subgoals''.  Completing only sequential subgoals yields full rewards (100 points), while each task under ``unordered and optional subgoals'' earns just 1 point. But many of these subgoals are important and necessary for achieving the final goals. Yet, they remain challenging even for strong models like GPT-3.5-Turbo. It is inappropriate to assign these necessary and difficult subgoals a score of 1. Moreover, the first subgoals in 'sequential subgoals' could be far from the initial state. This could result in very low rewards despite significant progress being made.
%\jh{is this 1-point subgoal important? You could argue it is important thus 1 point is a bad design, but ``it is difficult'' does not entail 1 point is inappropriate, I don't see the logic here} 
%To address these issues, we re-annotate subgoals to calculate $r_t^{subgoal}$, where
% Specifically, we incorporate necessary observations as part of the subgoals. 
%the rewards for these subgoals are more uniform and distributed evenly throughout the task. 
%To ensure that our annotated subgoals are necessary for achieving final goals, we restrict the use of tools and designated task completion rooms in the task descriptions. 
% We show more details of we annotated subgoals in the Appendix~\ref{sec: Re-annotation of Scienceworld}%\jh{plz elaborate this point further in the appendix}

%\jh{can you give examples on how the original sub-goals are re-annotated, and how we modify text to provide more information in the appendix?}

\textbf{\textit{Embodied} - BabyAI (BA)~\citep{chevalier2018babyai}}
is an interactive 20x20 grid environment where agents navigate and interact with objects within a limited sight range. 
The original setup uses image-based observations and tensor-based actions like “0: move left”. We adapted it to include a textual action space and descriptive textual observations.
% For each step, the environment returns a text description of the current observation, such as “There is a red ball 1 step to your right and 1 step ahead of you. There is a wall 2 steps ahead.” We also introduced high-level actions, such as "go to red ball 1” and “toggle and go through green locked door 1”, to expand the action space and enrich the semantic complexity of the environment. 
Furthermore, we re-annotate subgoals for progress rates to fix subgoal sparsity problem  in the original environment (Appendix \ref{sec: Annotation of babyai}).
% Unlike the previous reward score in BabyAI which awards a point only after a new object is found or pickup, requiring many steps to see progress in reward score, our new approach increases density of the rewards, requiring fewer steps to achieve them. 
% We re-annotated subgoals and calculate with the equation of $r_{t}^{\text{subgoal}}$. 
% Subgoals are re-annotated to update the progress rate whenever the agent makes progress, such as navigating to another room, finding a red ball, and picking it up in the problem “pickup a red ball”.


 


% \subsection{Game Environments}
%Evaluating LLM agents as strategic game playing agents demands strong planning ability of agents. We select two tasks requiring planning skills and strategic development.


\textbf{\textit{Game} - Jericho (JC)~\citep{hausknecht2020interactive}}
is a collection of text-based game environments staged in fictional worlds. This task is unique as it requires strong world modeling ability : agents could only gain information about the magic world through exploration and interaction. 
% For example, for the task that requires the agent to perform actions with magic, it cannot reason with pre-trained commonsense knowledge and must perform exploration to understand the rules of the magic world. 
The original games are too long (need 50-300 steps to finish for LLM agents with fixed context length. Therefore we rewrite the goal of each adventure to restrict the games to be finished within 15 subgoals. 
% For example, \emph{zork1} game requires the player to enter a dungeon and explore the dungeon to find a bar. We rewrite the goal as “You need to find your way into a secret passage where the entrance is in the living room of the house.” and the agent only needs to find the entrance to the dungeon, which can be finished in 8 steps. 
% We use the $r_t^{\text{subgoal}}$ as progress rate metrics, and we meticulously annotate the subgoals for each problem. 
% Each subgoal characterize that the agent has solved a small problem, e.g. “find the entrance to the house" $\rightarrow$ “enter the house" $\rightarrow$ “find the living room" $\rightarrow$ “discover a trap door" $\rightarrow$ “find the entrance to dungeon". 

\textbf{\textit{Game} - PDDL (PL)~\citep{vallati20152014}} is a set of strategic games defined with Planning Domain Definition Language (PDDL). 
We selected 4 representative games, \textit{Gripper, Barman, Blocksworld, Tyreworld} to benchmark LLM agents in diverse scenarios.
% , where the agent needs to move balls across rooms, make cocktails, rearrange blocks and pump up and install new tyres to cars. 
% This task is difficult as it requires multiple rounds of planned actions to finish a single subgoal and agents need to plan strategically to avoid repetitive steps. 
% For example, in \textit{Barman}, the player is given a menu, and is required to make a few cocktails with a few containers and ingredients. The agent could use a strategy of trying to use different containers each time to avoid repetitive cleaning and save steps. 
% While the commonly-used environment implementation \citep{silver2020pddlgym} requires the agent to interact with an environment with PDDL expressions, 
We adapted the environment implementation \citep{silver2020pddlgym} written in PDDL expressions to provide text-based observations for agents, enabling a natural language interface for LLM. We measure progress using a matching score, $r_t^{\text{match}}$, which assesses similarity between the current and goal states (Appendix \ref{sec: Annotation of pddl}).
% e.g. \texttt{clean-shaker(hand1, hand2, shaker)} and provides observations as set of predicates \texttt{ontable(shaker1) $\wedge$ empty(shaker1)}.  
% we write parser rules to offer a text-based observation to agents that allows LLMs to interact with natural language to be consistent with other tasks. 
% e.g.“Shaker1 is on the table. Shaker1 is empty” and enable the agents to interact with the environment with simple text commands, e.g. “clean-shaker shaker1 with hand1 while hand2 is empty.” 
% We curate 10-20 problems for each of the four domains by ourselves, ensuring the problems are multi-round and diverse. 
% We use the $r_t^{\text{match}}$ as progress rate metric. The matching score compares the similarity between the properties of current state and the goal state. 
% e.g. for the goal state ``Block a is on block b. Block b is on the table'', if at current state ``Block a is on the table. Block b is on the table'', then the matching score is 0.5. The agent will receive a 100\% progress rate only if all conditions of the goal state are satisfied.

% \subsection{Web-based Environments}
% Evaluating LLM's capability as a generalist agent in web-based scenarios has become pivotal~\citep{Shi2017WorldOB, deng2023mind2web}. 
%Web agent is expected to navigate the network efficiently and perform diverse tasks amidst highly dynamic, intricate, and multi-turn interactions. 
% Based on the task categorization, we've pinpointed two tasks of high recognition and quality: the specific network task, WebShop~\citep{yao2022webshop}, and the general network task, WebArena~\citep{zhou2023webarena}. The latter permits unrestricted access to any supported webpage.

\textbf{\textit{Web} - WebShop (WS)~\citep{yao2022webshop}}
is a network-based simulation environment for e-commerce experiences. %featuring a website with 1.18 million actual products, each with distinct labels and attributes. 
Based on the original implementation method~\citep{yao2022webshop,shinn2023reflexion}, we have improved the error feedback, including refining the observation for exceeding page limits and interacting with wrong objects. These enhancements contribute to the effective interaction of the LLM agent with the environment.
We also measure the distance of the current state to the final goal as the progress rate and expand the product scoring rules from~\citet{yao2022webshop} to derive the score (Appendix \ref{sec: Re-annotation of WebShop}). 
% In this environment, the agent is allowed to interact with the system through `search[QUERY]' or `click[ELEMENT]' actions to purchase products matching the instructions. 


% As there are no sub-goals in the environment, to obtain a continuous progress rate, we expanded the calculation rules from~\citep{yao2022webshop}, calculating the score at different web pages (stages). 
% To measure the distance of the current state to the final goal as the progress rate, we expanded the product scoring rules from~\citet{yao2022webshop} to derive the score at different web pages. Please refer to Appendix~\ref{sec: Re-annotation of WebShop} for details.
% Specifically, we calculate the highest score of products on a page for both search result page and product description page. At the confirm order page, we calculate the final product score (a full score cannot be achieved if attributes are not correctly chosen). In this rule-based evaluation, the final progress rate is the average of the scores from the three pages.

\textbf{\textit{Web} - WebArena (WA)~\citep{zhou2023webarena}}
is a real web environment featuring various scenarios including business content and discussion forums. %containing four applications: online shopping, discussion forums, collaborative development, and business content management. 
% It supports 11 different web browsing actions. 
% such as click (element),  new tab, goto (URL), etc., and offers additional tools like maps and wikis. 
% The observation space consists of structured web content. 
% (the accessibility tree.
% \footnote{https://developer.mozilla.org/en-US/docs/Glossary/Accessibility\_tree}). Completing tasks in this highly realistic environment requires the agent to possess strong memory, high-level planning, common sense, and reasoning abilities.
%Compared to other datasets~\citep{deng2023mind2web,shi2017world}, WebArena offers multi-round and continuous web browsing interaction simulation. 
% We filtered 245 instances from the original dataset for two main sub-tasks: Site Navigation and Contact \& Config, each annotated with the target URLs or required content. 
To obtain the progress rate, we revised the existing method for calculating the final score~\citep{zhou2023webarena} and continuously computed the progress rate at each step, fusing the URL matching score with the content matching score, 
% derived from the current URL and target URL, with the content matching score calculated based on the detected required content, 
as detailed in Appendix~\ref{sec: Re-annotation of WebArena}.
% The progress rate is formulated as follows: 

% \begin{equation}
%     r^{\text{match}} = \frac{n}{m+n}(r_{d}(r_{q}+r_{p})) + \frac{m}{n+m}r_{c} \quad (n=3; m=0,1,2,\ldots)
% \end{equation}

% In our method, we effectively utilize the annotation data, treating URLs as indicators of the web browsing trajectory and required contents as integral scoring points. First, We use \texttt{util.parse} to dissect URL into domain, query, and parameters. $r_{d}$ is assigned a binary value (0/1), scoring only if the domain is correct. Then, by applying the LCS (Longest Common Subsequence) algorithm, we calculate the matching score $r_{q}$ between the current query and the target query, and $r_{p}$ is the F1 score between the current and target parameters. The content matching score $r_{c}$ is calculated based on the proportion of required contents detected at each step. Ultimately, the progress rate is computed as a weighted sum of the URL matching and content matching scores. Here, $n$ represents the number of URL components, and $m$ denotes the count of required contents.

% \subsection{Tool Environments}
% LLM agents synergize with external tools; however, current benchmarks miss multi-round complexity or progress annotations~\citep{ qin2023toolllm, liu2023agentbench}. \BenchName{} features tool environments in two categories.
%~\citep{mialon2023augmented,qin2023tool,shen2023hugginggpt,liang2023taskmatrix}
% Another application area of LLM agents lies in their synergetic interactions with external tools.
% However, existing tool-using benchmarks either do not require complex multi-round interaction~\citep{li2023api} or lack annotated data for progress rate tracking~\citep{qin2023toolllm, liu2023agentbench}.
% %meticulously 
% In \BenchName{}, we have established five distinct tool environments that are categorized into two primary groups.
% Within these environment settings, human annotators are tasked with designing goals requiring complex multi-round interaction to accomplish.
% As indicated in Table~\ref{tab: Statistics of 10 environments}, these tool environments typically necessitate an average of five turns to complete.
% For detailed information about how we construct tool environments, please refer to Appendix~\ref{appendix:tool}.

%\textbf{Tool-Query~(TQ)}
%consists of three sub-environments: Weather Environment, Movie Environment and Academia Environment.
%Tools involved primarily serve the purpose of querying information.
% about weather, movie and computer science academia.
% In Tool-Query Environments, we employ $r_{t}^{\text{subgoal}}$ as a metric to measure progress rate.
% While designing actions for these environments, we ensure that each action's functionality is indecomposable~(i.e., the functionality and outcome of one action can not be achieved through other actions).
% This design choice results in a deterministic set of required golden actions to achieve our annotated goal.
% Furthermore,
%To compute the progress rate $r_{t}^{\text{subgoal}}$,
%the authors manually identify golden actions for each goal.
%Outputs returned by executing these golden actions are then processed as subgoals.
% For example, in Table~\ref{tab:goal_trajectory}, to accomplish the goal ``Avatar versus Forrest Gump, which has a higher rating?'', the action ``get\_movie\_details with Action Input: \{``movie id'': ``13''\}'' is one of the golden actions.
% We take the output of it as a subgoal, which contains information about the ``vote\_average'' of ``Forrest Gump''.
% \jh{I don't understand this subgoal description and definition, confusing}

%\textbf{Tool-Operation~(TO)}
%comprises two sub-environments: Todo Environment and Sheet Environment.
%Tools incorporated within these environments are designed for accessing and modifying information. 
% about personal agendas and spreadsheets.
%For Todo Environment, we adopt $r_{t}^{\text{subgoal}}$ as progress rate metric as Tool-Query Environments.
% Subgoals are annotated following the same methodology as Tool-Query Environments.
%In the Sheet Environment, progress rate is assessed using the $r_{t}^{\text{match}}$.
%The authors manually annotate the golden spreadsheet for each goal.
% After each interaction round, 
%We evaluate the matching score based on the proportion of cells in current spreadsheet that align with the golden spreadsheet.
% In Table~\ref{tab:goal_trajectory}, we provide an example of goal and trajectory in Sheet Environment.
\textbf{\textit{Tool} - Tool-Query (TQ)}
consists of three sub-environments: Weather, Movie
and Academia Environment. This tasks primarily involves querying information
from respective databases by planning the use of diverse query tools. We manually curate diverse problems for each environment. We also annotate subgoals to compute the progress rate $r_{t}^{\text{subgoal}}$ (Appendix \ref{sec: Annotation of tool query}). While LLMs often provide direct answers in question answering, we only consider an answer correct if the model follows the appropriate trajectory to access the databases. To ensure this, we design questions that cannot be answered directly by state-of-the-art LLMs and provide in-context examples to guide the LLM in querying the databases effectively.

\textbf{\textit{Tool} - Tool-Operation (TO)} includes two sub-environments: Todo list management and Google Sheet Operations. These tasks involve using tools to access and modify information. The progress rate in the Todo Environment is measured using $r_{t}^{\text{subgoal}}$, similar to Tool-Query Environments. In the Sheet Environment, progress is evaluated using $r_{t}^{\text{match}}$, which uses a matching score between the cells of the current and golden table (Appendix \ref{sec: Annotation of tool operation}).


\begin{figure*}[!t]%
    \centering
{\includegraphics[width=\linewidth]{figures/progress_rate_vs_human_evaluation_score.pdf} }%
\vspace{-5mm}
    \caption{Human verification of \BenchName{} progress rate. The Pearson correlation coefficients $\rho$ compare human evaluations with the progress rates on 60 different trajectories per task. The Fleiss' kappa $\kappa$ reflects inter-annotator agreement. The trajectories are generated by three models: GPT-4, GPT-3.5-Turbo, and DeepSeek-67b.
    % \jh{it seems every task should have 60 points in the figure right? But why some tasks have much fewer points?} \chang{due to overlapping of some points}
    }
    \label{fig:human study}%
    \vspace{-5mm}
\end{figure*}

\subsection{Annotation Verification and Metric Justification\label{sec: annotation verification}}


After human annotation, we manually verified our labeled subgoals through multiple verification stages to ensure its quality, as detailed in Appendix \ref{sec: Data Quality Control}. 
More importantly, we conduct a user study to justify our proposed progress rate metric, 
asking human annotators to assess the progress of model trajectories and then evaluating its correlation with our automatic progress rate metric. 
Specifically, we gather 60 model trajectories for each of the 8 tasks from three strong LLMs——\texttt{GPT-4}, \texttt{GPT-3.5-Turbo}, and \texttt{Deepseek-67b}. Each trajectory is assessed by four authors of the paper. The individual human rater is asked to select progress score from $\{0\%, 25\%, 50\%, 75\%, 100\%\}$ given trajectories and task descriptions, without seeing the automatic score.   
Mean of four scores is taken as the final human score for every trajectory. 
We show the Pearson correlation between human progress score and the progress rate in Figure~\ref{fig:human study}, and report Fleiss'kappa $\kappa$ to reflect inter-annotator agreement. 
Results show that progress rate highly correlates with human assessment on the progress where the Pearson correlation exceeds 0.95 on all tasks, 
and substantial agreement is reached among the annotators.
% \jh{I rewrite this a bit, this part is actually very important, thus I try to give more details to emphasize it.}

% assessed the correlation between human-scored trajectory progress and the \BenchName{} progress rate. We gathered 60 trajectories from three different LLMs, with each trajectory scored by four individual evaluators using scores from $\{0, 0.25, 0.5, 0.75, 1\}$. As Figure \ref{fig:human study} illustrates, there is a strong correlation, exceeding a Pearson correlation of 0.95 on all tasks, between human evaluations and progress rate metrics, confirming the reliability of our progress rate assessments.

% To ensure the quality of our labeled subgoals, we conduct manual validation on the annotated subgoals. 
% We design an interactive UI (Figure \ref{fig: sub-goal annotations 2}).We conduct at least two validation rounds. For each task, the first round is self-check, where the annotators play with all of their annotated tasks. In the second round, we asked two graduate students to sample environments and interact with them. During this interaction, annotators are provided with their progress rate. If the progress rate seemed unreasonable to the annotators, the task is recorded and re-annotated. 
% We report our error rate in the second round in Table~\ref{tab: ration of environments with annotation erros}. This indicates that our final annotation has a very low error rate and is of high quality. More details on our data check procedures are provided in Appendix \ref{sec: Data Quality Control}.


%\jh{we should briefly mention the annotation error is low here. This paragraph basically says we sample a subset and reannotate the errors -- this does not entail our final annotation is of high quality} 

% The matching score is computed by dividing the number of matching cells by the total number of cells.


% For Tool-Query Environments and Todo Environment,


% This ratio provides a fine-grained score representing the percentage of agreement between the sheets.
% because golden action paths are diverse, we cannot annotate a unique subgoal set~(e.g., an agent can ``transpose the spreadsheet'' then ``sort the first column of the spreadsheet'' or ``sort the first row of the spreadsheet'' then ``transpose the spreadsheet'', these two action paths cause the same final spreadsheet but different intermediate spreadsheet).
% Thus, we do not annotate subgoals for Sheet Environment.
% \jh{does this mean we do not have progress rate metric for this env?}
% \yc{We use \textbf{matching score}~(like Web tasks) as the progress rate metric for sheet environment}

% \begin{figure}[!t]%
%     \centering
%     {\includegraphics[width=1.0\linewidth]{figures/tool_subgoals.pdf} }%
%     \caption{The process of annotating subgoals for a goal in Todo Environment.}
%     \label{fig:tool_subgoals}%
% \end{figure}
